% --- Archon physics formalization source begin ---
% archon:physics
% archon:covers PhyXMiniProblems/problem_phyx_mini_0351.lean
% archon:source-report reports/phyx_mini/problem_phyx_mini_0351.source.json

\chapter{Physics problem phyx\_mini\_0351}
\label{ch:PhyXMiniProblems_problem_phyx_mini_0351}

\paragraph{Problem source.}
\#\# Physical scenario

A heat engine operates using the cycle shown in figure. The working substance is 2.00 mol of helium gas, which reaches a maximum temperature of \$327\textasciicircum{}\{\textbackslash{}circ\}C\$. Assume the helium can be treated as an ideal gas. Process \$bc\$ is isothermal. The pressure in states \$a\$ and \$c\$ is \$1.00 \textbackslash{}times 10\textasciicircum{}5\$ Pa, and the pressure in state \$b\$ is \$3.00 \textbackslash{}times 10\textasciicircum{}5\$ Pa.

\#\# Question

What is its efficiency?

\#\# Answer choices

- A: A: 10\textbackslash{}\%

- B: B: 8.7\textbackslash{}\%

- C: C: 15.9\textbackslash{}\%

- D: D: 12.3\textbackslash{}\%

\#\# Figure caption (auxiliary; use the image as primary evidence)

The image features a graph with a set of labeled points and lines forming a closed loop. The axes are labeled with "O" on the vertical axis and "V" on the horizontal axis.

- There are three labeled points: \textbackslash{}( a \textbackslash{}), \textbackslash{}( b \textbackslash{}), and \textbackslash{}( c \textbackslash{}).
- The points are connected by blue arrows forming a triangular path:
  - An arrow goes vertically upwards from \textbackslash{}( a \textbackslash{}) to \textbackslash{}( b \textbackslash{}).
  - A curved arrow goes from \textbackslash{}( b \textbackslash{}) to \textbackslash{}( c \textbackslash{}).
  - A horizontal arrow goes from \textbackslash{}( c \textbackslash{}) back to \textbackslash{}( a \textbackslash{}).

Each arrow indicates direction, suggesting cyclical movement between the points. The style is clean, with black points and blue arrows forming the shape.

\#\# Dataset metadata

Laws of Thermodynamics; Physical Model Grounding Reasoning, Multi-Formula Reasoning

\paragraph{Recorded answer/context.}
B: B: 8.7\textbackslash{}\%

\paragraph{Figure/image path.}
/root/proposal\_for\_physic/science-mango/phyx\_mini\_run/phyx\_data/test\_image/351.png

\paragraph{Formalization target.}
create a compiling Lean file with sorry bodies at `PhyXMiniProblems/problem\_phyx\_mini\_0351.lean`. The Lean declarations must preserve the physical quantities, dimensions or dimensional roles, figure labels, governing-law hypotheses, and final relation expressed by this problem.
Use Mathlib/Physlib names found through LeanExplore where available. If a physics API is missing, introduce faithful local abstractions rather than scalar placeholder aliases.

\begin{theorem}[Physics formalization target]
\label{thm:physics:phyx_mini_0351:target}
\lean{PhyXMiniProblems.ProblemPhyXMini0351.problem_phyx_mini_0351}
\uses{def:physics:phyx-mini-0351:phyxminiproblems-problemphyxmini0351-heliumheatenginecycle, def:physics:phyx-mini-0351:phyxminiproblems-problemphyxmini0351-matchesheatenginescenario, def:physics:phyx-mini-0351:phyxminiproblems-problemphyxmini0351-matchesproblemreadouts, def:physics:phyx-mini-0351:phyxminiproblems-problemphyxmini0351-hasphysicalheatengineparameters, def:physics:phyx-mini-0351:phyxminiproblems-problemphyxmini0351-matchessuppliedpvcyclefigure, def:physics:phyx-mini-0351:phyxminiproblems-problemphyxmini0351-satisfiesidealmonatomicgascyclelaws, def:physics:phyx-mini-0351:phyxminiproblems-problemphyxmini0351-answerchoice, def:physics:phyx-mini-0351:phyxminiproblems-problemphyxmini0351-efficiencyroundstopercent, def:physics:phyx-mini-0351:phyxminiproblems-problemphyxmini0351-efficiencyroundstochoice}
The assigned autoformalize agent should translate this physics problem into Lean declarations in the covered file, with theorem and lemma proof bodies written as `by sorry`.
\end{theorem}
\begin{proof}
This is an autoformalization task, not a proof task. Produce statements that can later be proved by the physics prover without weakening the physical model.
\end{proof}

% --- Archon declaration topology begin ---
\section{Declaration topology}
\begin{definition}[Gas Volume]
  \label{def:physics:phyx-mini-0351:phyxminiproblems-problemphyxmini0351-gasvolume}
  \lean{PhyXMiniProblems.ProblemPhyXMini0351.GasVolume}
  A nonnegative physical volume, independent of the selected unit system
\end{definition}
\begin{proof}
  This is the declared model definition of Gas Volume.
\end{proof}
\begin{definition}[pressure In Pascals]
  \label{def:physics:phyx-mini-0351:phyxminiproblems-problemphyxmini0351-pressureinpascals}
  \lean{PhyXMiniProblems.ProblemPhyXMini0351.pressureInPascals}
  Read a physical pressure in pascals
\end{definition}
\begin{proof}
  This is the declared model definition of pressure In Pascals.
\end{proof}
\begin{definition}[volume In Cubic Meters]
  \label{def:physics:phyx-mini-0351:phyxminiproblems-problemphyxmini0351-volumeincubicmeters}
  \lean{PhyXMiniProblems.ProblemPhyXMini0351.volumeInCubicMeters}
  \uses{def:physics:phyx-mini-0351:phyxminiproblems-problemphyxmini0351-gasvolume}
  Read a physical volume in cubic metres
\end{definition}
\begin{proof}
  This is the declared model definition of volume In Cubic Meters.
\end{proof}
\begin{definition}[energy In Joules]
  \label{def:physics:phyx-mini-0351:phyxminiproblems-problemphyxmini0351-energyinjoules}
  \lean{PhyXMiniProblems.ProblemPhyXMini0351.energyInJoules}
  Read a signed physical energy in joules
\end{definition}
\begin{proof}
  This is the declared model definition of energy In Joules.
\end{proof}
\begin{definition}[temperature In Kelvins]
  \label{def:physics:phyx-mini-0351:phyxminiproblems-problemphyxmini0351-temperatureinkelvins}
  \lean{PhyXMiniProblems.ProblemPhyXMini0351.temperatureInKelvins}
  Read a Physlib absolute temperature in kelvins for this SI model
\end{definition}
\begin{proof}
  This is the declared model definition of temperature In Kelvins.
\end{proof}
\begin{definition}[Thermodynamic State]
  \label{def:physics:phyx-mini-0351:phyxminiproblems-problemphyxmini0351-thermodynamicstate}
  \lean{PhyXMiniProblems.ProblemPhyXMini0351.ThermodynamicState}
  The three labelled thermodynamic states in the supplied figure
\end{definition}
\begin{proof}
  This is the declared model definition of Thermodynamic State.
\end{proof}
\begin{definition}[Cycle Leg]
  \label{def:physics:phyx-mini-0351:phyxminiproblems-problemphyxmini0351-cycleleg}
  \lean{PhyXMiniProblems.ProblemPhyXMini0351.CycleLeg}
  The three directed legs of the cycle
\end{definition}
\begin{proof}
  This is the declared model definition of Cycle Leg.
\end{proof}
\begin{definition}[leg Start]
  \label{def:physics:phyx-mini-0351:phyxminiproblems-problemphyxmini0351-legstart}
  \lean{PhyXMiniProblems.ProblemPhyXMini0351.legStart}
  \uses{def:physics:phyx-mini-0351:phyxminiproblems-problemphyxmini0351-thermodynamicstate, def:physics:phyx-mini-0351:phyxminiproblems-problemphyxmini0351-cycleleg}
  Initial state of each directed cycle leg
\end{definition}
\begin{proof}
  This is the declared model definition of leg Start.
\end{proof}
\begin{definition}[leg Finish]
  \label{def:physics:phyx-mini-0351:phyxminiproblems-problemphyxmini0351-legfinish}
  \lean{PhyXMiniProblems.ProblemPhyXMini0351.legFinish}
  \uses{def:physics:phyx-mini-0351:phyxminiproblems-problemphyxmini0351-thermodynamicstate, def:physics:phyx-mini-0351:phyxminiproblems-problemphyxmini0351-cycleleg}
  Final state of each directed cycle leg
\end{definition}
\begin{proof}
  This is the declared model definition of leg Finish.
\end{proof}
\begin{definition}[Thermodynamic Process Kind]
  \label{def:physics:phyx-mini-0351:phyxminiproblems-problemphyxmini0351-thermodynamicprocesskind}
  \lean{PhyXMiniProblems.ProblemPhyXMini0351.ThermodynamicProcessKind}
  Thermodynamic process types used in the problem
\end{definition}
\begin{proof}
  This is the declared model definition of Thermodynamic Process Kind.
\end{proof}
\begin{definition}[Gas Species]
  \label{def:physics:phyx-mini-0351:phyxminiproblems-problemphyxmini0351-gasspecies}
  \lean{PhyXMiniProblems.ProblemPhyXMini0351.GasSpecies}
  The working-gas species stated in the problem
\end{definition}
\begin{proof}
  This is the declared model definition of Gas Species.
\end{proof}
\begin{definition}[Figure State Position]
  \label{def:physics:phyx-mini-0351:phyxminiproblems-problemphyxmini0351-figurestateposition}
  \lean{PhyXMiniProblems.ProblemPhyXMini0351.FigureStatePosition}
  Qualitative locations of state markers in the raster diagram
\end{definition}
\begin{proof}
  This is the declared model definition of Figure State Position.
\end{proof}
\begin{definition}[Figure Path Shape]
  \label{def:physics:phyx-mini-0351:phyxminiproblems-problemphyxmini0351-figurepathshape}
  \lean{PhyXMiniProblems.ProblemPhyXMini0351.FigurePathShape}
  Shapes of the three blue directed paths in the raster diagram
\end{definition}
\begin{proof}
  This is the declared model definition of Figure Path Shape.
\end{proof}
\begin{definition}[Supplied PVCycle Figure]
  \label{def:physics:phyx-mini-0351:phyxminiproblems-problemphyxmini0351-suppliedpvcyclefigure}
  \lean{PhyXMiniProblems.ProblemPhyXMini0351.SuppliedPVCycleFigure}
  \uses{def:physics:phyx-mini-0351:phyxminiproblems-problemphyxmini0351-thermodynamicstate, def:physics:phyx-mini-0351:phyxminiproblems-problemphyxmini0351-cycleleg, def:physics:phyx-mini-0351:phyxminiproblems-problemphyxmini0351-figurestateposition, def:physics:phyx-mini-0351:phyxminiproblems-problemphyxmini0351-figurepathshape}
  Literal labels, marker locations, and directed paths in the supplied image
\end{definition}
\begin{proof}
  This is the declared model definition of Supplied PVCycle Figure.
\end{proof}
\begin{definition}[Helium Heat Engine Cycle]
  \label{def:physics:phyx-mini-0351:phyxminiproblems-problemphyxmini0351-heliumheatenginecycle}
  \lean{PhyXMiniProblems.ProblemPhyXMini0351.HeliumHeatEngineCycle}
  \uses{def:physics:phyx-mini-0351:phyxminiproblems-problemphyxmini0351-gasvolume, def:physics:phyx-mini-0351:phyxminiproblems-problemphyxmini0351-thermodynamicstate, def:physics:phyx-mini-0351:phyxminiproblems-problemphyxmini0351-cycleleg, def:physics:phyx-mini-0351:phyxminiproblems-problemphyxmini0351-thermodynamicprocesskind, def:physics:phyx-mini-0351:phyxminiproblems-problemphyxmini0351-gasspecies, def:physics:phyx-mini-0351:phyxminiproblems-problemphyxmini0351-suppliedpvcyclefigure}
  The physical heat engine. Heat is positive into the gas and work is positive when done by the gas. Consequently both are signed DimEnergy quantities. molarGasConstantJoulesPerMoleKelvin is an explicit coherent-SI readout, and thermalEfficiency is dimensionless
\end{definition}
\begin{proof}
  This is the declared model definition of Helium Heat Engine Cycle.
\end{proof}
\begin{definition}[Matches Heat Engine Scenario]
  \label{def:physics:phyx-mini-0351:phyxminiproblems-problemphyxmini0351-matchesheatenginescenario}
  \lean{PhyXMiniProblems.ProblemPhyXMini0351.MatchesHeatEngineScenario}
  \uses{def:physics:phyx-mini-0351:phyxminiproblems-problemphyxmini0351-heliumheatenginecycle}
  The gas identity and the thermodynamic interpretation of the three legs
\end{definition}
\begin{proof}
  This is the declared model definition of Matches Heat Engine Scenario.
\end{proof}
\begin{definition}[Matches Problem Readouts]
  \label{def:physics:phyx-mini-0351:phyxminiproblems-problemphyxmini0351-matchesproblemreadouts}
  \lean{PhyXMiniProblems.ProblemPhyXMini0351.MatchesProblemReadouts}
  \uses{def:physics:phyx-mini-0351:phyxminiproblems-problemphyxmini0351-pressureinpascals, def:physics:phyx-mini-0351:phyxminiproblems-problemphyxmini0351-temperatureinkelvins, def:physics:phyx-mini-0351:phyxminiproblems-problemphyxmini0351-heliumheatenginecycle}
  The numerical information printed in the problem statement
\end{definition}
\begin{proof}
  This is the declared model definition of Matches Problem Readouts.
\end{proof}
\begin{definition}[Has Physical Heat Engine Parameters]
  \label{def:physics:phyx-mini-0351:phyxminiproblems-problemphyxmini0351-hasphysicalheatengineparameters}
  \lean{PhyXMiniProblems.ProblemPhyXMini0351.HasPhysicalHeatEngineParameters}
  \uses{def:physics:phyx-mini-0351:phyxminiproblems-problemphyxmini0351-pressureinpascals, def:physics:phyx-mini-0351:phyxminiproblems-problemphyxmini0351-volumeincubicmeters, def:physics:phyx-mini-0351:phyxminiproblems-problemphyxmini0351-energyinjoules, def:physics:phyx-mini-0351:phyxminiproblems-problemphyxmini0351-temperatureinkelvins, def:physics:phyx-mini-0351:phyxminiproblems-problemphyxmini0351-heliumheatenginecycle}
  Positivity assumptions selecting the physical branch of the model
\end{definition}
\begin{proof}
  This is the declared model definition of Has Physical Heat Engine Parameters.
\end{proof}
\begin{definition}[Matches Supplied PVCycle Figure]
  \label{def:physics:phyx-mini-0351:phyxminiproblems-problemphyxmini0351-matchessuppliedpvcyclefigure}
  \lean{PhyXMiniProblems.ProblemPhyXMini0351.MatchesSuppliedPVCycleFigure}
  \uses{def:physics:phyx-mini-0351:phyxminiproblems-problemphyxmini0351-legstart, def:physics:phyx-mini-0351:phyxminiproblems-problemphyxmini0351-legfinish, def:physics:phyx-mini-0351:phyxminiproblems-problemphyxmini0351-heliumheatenginecycle}
  Data transcribed directly from the primary p--V raster image
\end{definition}
\begin{proof}
  This is the declared model definition of Matches Supplied PVCycle Figure.
\end{proof}
\begin{definition}[Satisfies Ideal Monatomic Gas Cycle Laws]
  \label{def:physics:phyx-mini-0351:phyxminiproblems-problemphyxmini0351-satisfiesidealmonatomicgascyclelaws}
  \lean{PhyXMiniProblems.ProblemPhyXMini0351.SatisfiesIdealMonatomicGasCycleLaws}
  \uses{def:physics:phyx-mini-0351:phyxminiproblems-problemphyxmini0351-pressureinpascals, def:physics:phyx-mini-0351:phyxminiproblems-problemphyxmini0351-volumeincubicmeters, def:physics:phyx-mini-0351:phyxminiproblems-problemphyxmini0351-energyinjoules, def:physics:phyx-mini-0351:phyxminiproblems-problemphyxmini0351-temperatureinkelvins, def:physics:phyx-mini-0351:phyxminiproblems-problemphyxmini0351-legstart, def:physics:phyx-mini-0351:phyxminiproblems-problemphyxmini0351-legfinish, def:physics:phyx-mini-0351:phyxminiproblems-problemphyxmini0351-heliumheatenginecycle}
  The governing laws are independent of the requested numerical efficiency: P V = n R T at every state, in coherent SI readouts; U = (3/2) n R T for monatomic helium; the first law Delta U = Q - W on every directed leg; zero isochoric work, logarithmic isothermal work, and P Delta V isobaric work; cycle work and positive-heat-input bookkeeping; and the definition eta = W\_net / Q\_in of heat-engine efficiency
\end{definition}
\begin{proof}
  This is the declared model definition of Satisfies Ideal Monatomic Gas Cycle Laws.
\end{proof}
\begin{definition}[Answer Choice]
  \label{def:physics:phyx-mini-0351:phyxminiproblems-problemphyxmini0351-answerchoice}
  \lean{PhyXMiniProblems.ProblemPhyXMini0351.AnswerChoice}
  The answer letters shown in the multiple-choice problem
\end{definition}
\begin{proof}
  This is the declared model definition of Answer Choice.
\end{proof}
\begin{definition}[answer Efficiency Percent]
  \label{def:physics:phyx-mini-0351:phyxminiproblems-problemphyxmini0351-answerefficiencypercent}
  \lean{PhyXMiniProblems.ProblemPhyXMini0351.answerEfficiencyPercent}
  \uses{def:physics:phyx-mini-0351:phyxminiproblems-problemphyxmini0351-answerchoice}
  Displayed efficiencies, read as percentages
\end{definition}
\begin{proof}
  This is the declared model definition of answer Efficiency Percent.
\end{proof}
\begin{definition}[recorded Dataset Answer]
  \label{def:physics:phyx-mini-0351:phyxminiproblems-problemphyxmini0351-recordeddatasetanswer}
  \lean{PhyXMiniProblems.ProblemPhyXMini0351.recordedDatasetAnswer}
  \uses{def:physics:phyx-mini-0351:phyxminiproblems-problemphyxmini0351-answerchoice}
  The answer label recorded by the source dataset; it is never a premise
\end{definition}
\begin{proof}
  This is the declared model definition of recorded Dataset Answer.
\end{proof}
\begin{definition}[Efficiency Rounds To Percent]
  \label{def:physics:phyx-mini-0351:phyxminiproblems-problemphyxmini0351-efficiencyroundstopercent}
  \lean{PhyXMiniProblems.ProblemPhyXMini0351.EfficiencyRoundsToPercent}
  \uses{def:physics:phyx-mini-0351:phyxminiproblems-problemphyxmini0351-heliumheatenginecycle}
  Agreement after rounding an efficiency to the nearest tenth of a percent
\end{definition}
\begin{proof}
  This is the declared model definition of Efficiency Rounds To Percent.
\end{proof}
\begin{definition}[Efficiency Rounds To Choice]
  \label{def:physics:phyx-mini-0351:phyxminiproblems-problemphyxmini0351-efficiencyroundstochoice}
  \lean{PhyXMiniProblems.ProblemPhyXMini0351.EfficiencyRoundsToChoice}
  \uses{def:physics:phyx-mini-0351:phyxminiproblems-problemphyxmini0351-heliumheatenginecycle, def:physics:phyx-mini-0351:phyxminiproblems-problemphyxmini0351-answerchoice, def:physics:phyx-mini-0351:phyxminiproblems-problemphyxmini0351-answerefficiencypercent, def:physics:phyx-mini-0351:phyxminiproblems-problemphyxmini0351-efficiencyroundstopercent}
  Agreement with the percentage printed beside an answer choice
\end{definition}
\begin{proof}
  This is the declared model definition of Efficiency Rounds To Choice.
\end{proof}
% --- Archon declaration topology end ---
% --- Archon physics formalization source end ---
